%! Inverse Functions — Unit 1, Lesson 1.6 — Warm-Up (Answer Key)
\documentclass[10pt]{article}
\usepackage{precalculus-article}
\usepackage{precalculus-key}   % pulls in -boxes; do NOT also load -boxes
\newcommand{\TallMath}[1]{$\displaystyle #1\rule[-1.4em]{0pt}{3.2em}$}

\begin{document}

\pageheader{Unit 1, Lesson 1.6}{Warm-Up --- Answer Key}

\vspace{0.12in}

\begin{enumerate}[itemsep=12pt]

\item Solve each equation for $x$. (Plain algebra today, load-bearing later.)

\begin{work}
      4x - 1 &= 11 \\
         4x &= 12 \\
          x &= 3 \\
  1.8x + 32 &= 68 \\
       1.8x &= 36 \\
          x &= 20
\end{work}

\begin{enumerate}[label=(\alph*), itemsep=4pt]
  \item $4x - 1 = 11$ \qquad $x = $ \ans{3}
  \item $1.8x + 32 = 68$ \qquad $x = $ \ans{20}
\end{enumerate}

\item From Lesson 1.5: let $f(x) = 4x - 1$ and $g(x) = \dfrac{x + 1}{4}$.
      Evaluate each composition.

\begin{work}
     g(7) &= \frac{7 + 1}{4} \\
          &= 2 \\
  f(g(7)) &= f(2) \\
          &= 7 \\
     f(7) &= 4(7) - 1 \\
          &= 27 \\
  g(f(7)) &= g(27) \\
          &= 7
\end{work}

$f(g(7)) = $ \ans{7} \qquad $g(f(7)) = $ \ans{7}

\item The table gives every value of a function $h$. Fill in the second table,
      which lists the \textbf{same pairs with the inputs and outputs traded}.

\smallskip
\renewcommand{\arraystretch}{1.4}
\begin{center}
\begin{tabular}{|c|c|c|c|c|}
\hline
$x$ & 1 & 2 & 3 & 4 \\
\hline
$h(x)$ & 5 & 8 & 11 & 14 \\
\hline
\end{tabular}
\qquad
\begin{tabular}{|c|c|c|c|c|}
\hline
$x$ & 5 & 8 & 11 & 14 \\
\hline
new output & \ans{1} & \ans{2} & \ans{3} & \ans{4} \\
\hline
\end{tabular}
\end{center}

\end{enumerate}

\end{document}
